\section{Discussion and conclusion}\label{sec:discussion}

\subsection{Discussion}\label{sec:08-discussion}

\paragraph{A structural ceiling, now measured.}
By \cref{prop:upper-bound} every predicted $\theta_k$ lies above the discrete eigenvalue $\lambda_{h,k}$ of the space in which the model runs, so a trained model cannot be more accurate than the finite-element solution of that space; on the six test geometries that eigenvalue lies at most \SI{0.047}{\percent} above the high-fidelity (HF) frequency (\SI{0.033}{\percent} on four earlier ones). % src: docs/34 §8.1; docs/34 §4.1; EDITOR_NOTES_2 dec. 9
Meshing and operator assembly are needed for every geometry, and at the sizes studied (about $7\times10^3$ tetrahedra on average) plain finite elements dominate the model (\cref{sec:experiments}): % src: docs/35 §3b (6969 tets mean)
on one thread, a first-order solve on the model mesh took \SI{1.7}{s} for a mean frequency error of \SI{1.25}{\percent}, a second-order solve on a mesh with 48\,\% fewer tetrahedra \SI{5.7}{s} for \SI{0.037}{\percent}, and the second-order solve on the model mesh \SI{13.9}{s} for \SI{0.019}{\percent}, against \SI{8.3}{s} and \SIrange{1.4}{1.7}{\percent} for the model path. % src: docs/35 §3b; docs/34 §8.2 (8.3 s, one geometry, not re-measured)
The model therefore has no speed or accuracy advantage over a direct solve here. % src: docs/35 §3b
Its accuracy is also far from the $10^{-5}$--$10^{-4}$ that final designs require. % src: docs/33 §6

\paragraph{A parametric baseline on one family.}
A Gaussian process on the six geometry parameters matched the model once 24 or more geometries were available: with 40 LF geometries \SI{1.01}{\percent} mean and \SI{1.82}{\percent} maximum frequency error against \SI{1.37}{\percent} and \SI{5.30}{\percent} for the model, with 24 geometries \SI{1.64}{\percent} against \SI{1.55}{\percent}. % src: docs/35 §3a
With 8 geometries the model was about 2.3 times better (\SI{1.66}{\percent} against \SI{3.76}{\percent}), and the GP was equally accurate with LF and HF labels. % src: docs/35 §3a
The GP returns no fields, and the comparison covers one family and single model runs.
At this scale the number of training geometries, not the label fidelity, governed the accuracy of both.

\paragraph{What is guaranteed.}
The bound is exact in exact arithmetic under the kernel assumption (A), which we verified in all 15 cases tested (five curved meshes, $p=1,2,3$, small meshes only). % src: docs/35 §1
For the implemented conjugate-gradient projection it holds up to a relative perturbation $\bar\eta\approx(1\text{--}3)\,\tau_{\mathrm{CG}}$, about $10^{-8}$ at the shipped $\tau_{\mathrm{CG}}=10^{-8}$, and the a~posteriori form $\tilde\theta_k\ge\lambda_{h,k}/(1+\bar\eta)$ held in every run with $\bar\eta<1$. % src: docs/35 §2
Violations of the plain bound arose only for bases with a large gradient admixture, by at most $1.5\times10^{-8}$ at the shipped tolerance ($2\times10^{-11}$ and $1.4\times10^{-10}$ at a gradient fraction of $0.5$, close to that of a trained basis). % src: docs/35 §2 (reading)
The drop rule for near-gradient directions protects only for $\tau_{\mathrm{CG}}\lesssim10^{-6}$ (a pure-gradient basis at $10^{-4}$ gave $\bar\eta=1$, $\theta\approx0$), and without the projection gradient-rich bases do give spurious modes. % src: docs/35 §2
It is an index-wise bound on discrete eigenvalues, not a mode identity or an accuracy bound.

\paragraph{Low-fidelity labels, and where the approach could pay off.}
The Ritz layer reaches exactly the eigenvalues of the low-fidelity (LF) solve, which makes LF labels natural (\cref{sec:multifidelity}); here LF and HF labels differ by at most \SI{0.047}{\percent} (frequency) and \SI{0.32}{\percent} (field), one to two orders of magnitude less than the model error, so the study does not probe the regime where label type matters. % src: docs/34 §8.1; referee §E
They do differ on wall-dominated quantities (\cref{sec:multifidelity}), so HF labels remain necessary for such losses and for test sets.
We have not shown an advantage over parametric regression on one low-dimensional family or over direct finite elements at about $10^4$ tetrahedra. We hypothesise, without evidence yet, that it may pay off for geometry sets without a shared low-dimensional parametrisation (arbitrary CAD), for predicted fields and figures of merit with a guaranteed eigenvalue ordering, for warm starts of large finite-element solves, and for larger meshes and GPUs.

\subsection{Limitations}\label{sec:08-limitations}

\begin{itemize}
\item \textbf{Scale and scope.} Small CPU runs: at most 40 training geometries, a 64-dimensional two-block trunk, 600 steps per arm, one family (pipe-terminated pillbox cavities) with a minimum fillet of 0.12 that rounds the geometry considerably, six held-out geometries. Absolute errors (\SIrange{1}{2}{\percent}) reflect an under-trained model. % src: BRIEF §1; docs/34 §5, §8.2; EDITOR_NOTES_2 dec. 10
\item \textbf{Confounded, single-run LF/HF study.} Each arm of \cref{tab:equal-cost} is one run without seeds; fidelity and number of geometries change together (LF 8 against HF 8 was not run). Fine-tuning on HF was inconclusive (its control run was lost), and two mixed LF+HF runs failed. HF is one third-order solve on the start mesh, so the adaptive path and the label/model-mesh split are not exercised. % src: docs/34 §3, §4.3, §8.2; EDITOR_NOTES_2 dec. 2, 3
\item \textbf{Baselines and checks.} The baseline comparisons use earlier single model runs (no checkpoint was re-evaluated) and a model time measured on another geometry; the assumption check covers meshes of up to about $6\times10^3$ free degrees of freedom and the stress test two geometries with synthetic bases. % src: docs/35 §1--3 (limitations)
\item \textbf{Not trained or not tested.} The figures-of-merit loss is disabled for the field model; in a cubic space or on a finer mesh the bound holds but the frequency error grows by 0.2 to 0.7 percentage points; deformed geometries (label projection error unmeasured) and the GPU path were not tested. % src: docs/31 §2, §3.2, §3.8; docs/34 §4.2; EDITOR_NOTES_2 dec. 6
\end{itemize}

\subsection{Future work}\label{sec:08-future}

\paragraph{Hybrid inference.}
The predicted subspace can start the LOBPCG eigensolver~\cite{knyazev2001toward} on the target space, so that accuracy is guaranteed by the solver.
Our label solver gives a first indication at about $3.6\times10^5$ degrees of freedom (copper pillbox, third order): a warm start from the previous adaptive level reached an error of $1.9\times10^{-8}$ in 10 iterations (\SI{36}{s}), a cold start $1.7\times10^{-8}$ in 40 iterations (\SI{123}{s}). % src: docs/31 §3.4 (358k DOF)
With the \SI{26}{s} set-up and the \SI{68}{s} mode transfer the totals are about \SI{130}{s} (warm), \SI{149}{s} (cold) and \SI{175}{s} (direct), a marginal gain and a property of the label solver, not a model result. % src: docs/31 §3.4: 36+26+68 = 130 s; 123+26 = 149 s; direct 175 s; EDITOR_NOTES_2 dec. 12
Whether a model subspace is a good start is untested.

\paragraph{Training, families and validation.}
Next steps: the LF-8 versus HF-8 control with seeds, scaled LF+HF training with fidelity-weight sweeps, a wall-quantity loss, and families such as quarter-wave resonators, drift-tube cells, RF-dipole cavities and guns, where the comparison with parametric regression is the decisive test. % src: docs/34 §6, §8.3; docs/33 §4, §8; docs/31 §5, §6
A subset of each family should be cross-checked with an independent finite-element code such as Palace~\cite{awslabs2023palace}, and GPU measurements are needed.

\subsection{Conclusion}\label{sec:08-conclusion}

The surrogate predicts a basis of fields from which a conforming-space projection and a Rayleigh--Ritz step produce eigenpairs that are upper bounds of the discrete eigenvalues of the space they run in, exactly in exact arithmetic and up to an a~posteriori error of order $\tau_{\mathrm{CG}}$ for the implemented projection, without spurious gradient modes.
A single second-order solve on the model mesh is a cheap label: 2.8 to 4.3 times cheaper than a third-order one, depending on load, and at most \SI{0.047}{\percent} from it in frequency; % src: docs/34 §8.1; EDITOR_NOTES_2 dec. 4
models trained on such labels were comparable to one trained on HF labels (\cref{tab:equal-cost}), with a LF$\to$HF gap of 0.02 percentage points. % src: docs/34 §8.2
At this scale the number of geometries mattered more than the label fidelity, and we found no accuracy or speed advantage over parametric regression on one family or over direct finite elements at these sizes.
The contribution is the construction with its discrete guarantee, the data pipeline and controlled small-scale evidence; whether it pays off on arbitrary CAD geometries, at larger sizes or with warm starts remains to be tested.
